\documentclass[12pt,a4paper]{amsart}

\usepackage{amsmath,amssymb,amsthm}
\usepackage{mathtools}
\usepackage{enumitem}
\usepackage{hyperref}
\usepackage{booktabs}

%%%─── Theorem environments
\newtheorem{theorem}{Theorem}[section]
\newtheorem{lemma}[theorem]{Lemma}
\newtheorem{proposition}[theorem]{Proposition}
\newtheorem{corollary}[theorem]{Corollary}
\newtheorem{conjecture}[theorem]{Conjecture}
\theoremstyle{definition}
\newtheorem{definition}[theorem]{Definition}
\newtheorem{problem}[theorem]{Open Problem}
\theoremstyle{remark}
\newtheorem{remark}[theorem]{Remark}
\newtheorem*{remark*}{Remark}

%%%─── Macros
\newcommand{\Phihat}{\widehat{\Phi}}
\newcommand{\Ai}{\mathrm{Ai}}
\newcommand{\norm}[1]{\|#1\|}
\newcommand{\ip}[2]{\langle #1,\, #2\rangle}
\newcommand{\Hc}{\widetilde{\mathcal{H}}_c}
\newcommand{\Hlim}{\mathcal{H}_{\lim}}
\newcommand{\Hstrong}{H_{\lim}^{(\mathrm{str})}}
\newcommand{\Hspectral}{H_{\lim}^{(\mathrm{sp})}}
\newcommand{\HSOT}{\mathcal{H}^{\mathrm{SOT}}}
\newcommand{\Dom}{\mathrm{Dom}}
\newcommand{\qc}{q_c}
\newcommand{\qlim}{q_{\lim}}
\newcommand{\Fc}{\mathcal{F}}
\newcommand{\Bc}{B_c}
\DeclareMathOperator{\sgn}{sgn}

%%%─── AMS metadata
\subjclass[2020]{47A07, 47B25, 47A10, 42C05, 11M06}
\keywords{Quadratic forms, Mosco convergence,
resolvent convergence,
prolate spheroidal wave functions,
Riemann zeta function}

\begin{document}

\title[Mosco Convergence and Resolvent Stability
  of PSWF Concentration Forms]{%
  Mosco Convergence, Resolvent Stability,
  and the Bridge Theorem for
  PSWF Concentration Operators}

\author{[Author]}
\address{[Address]}
\email{[Email]}

\maketitle

\begin{abstract}
Let $\Hc:=P^*\mathcal{H}_c^{[0,1]}P$
be the Slepian--Pollak concentration
operator on $L^2(\mathbb{R})$.
The identity
$\ip{\Hc f}{f}=\|\Bc(P^*Pf)\|^2$
(coordinatisation lemma)
reduces $\Hc$ to a conjugated Fourier
projection; all results follow from
this identity and
$[-c,c]\subset[-c',c']$.
Two methodologically independent
threads --- (A) uniform bound $+$
Cauchy nets $\Rightarrow$
$\Hc\xrightarrow{\mathrm{SOT}}\Hstrong$;
(B) monotonicity $+$ compactness
$\Rightarrow$
$\qc\xrightarrow{\mathrm{M}}\qlim$ ---
are unified by $T_{\qlim}=\Hstrong$,
yielding strong resolvent convergence.
For a spectrally constructed
$\Hspectral$ (Paper~IX),
the identification
$\Hstrong=\overline{\Hspectral}$
(Bridge theorem) is conditional
on a structural hypothesis
(Paper~IX, Open Problem~7);
see Sec.~\ref{sec:bridge}.
No unconditional connection to
$\zeta(s)$-zeros is made.
\end{abstract}

\section*{Guide for the Referee}

\[
\boxed{\text{FTB}}
\to
\begin{cases}
  \text{Mono.}+\text{HS}
  \Rightarrow
  \qc\xrightarrow{\mathrm{M}}\qlim
  &\text{(Thread B)}\\
  \text{Bound}+\text{density}
  \Rightarrow
  \Hc\xrightarrow{\mathrm{SOT}}\Hstrong
  &\text{(Thread A)}
\end{cases}
\xRightarrow{T_{\qlim}=\Hstrong}
\text{resolvent}
\xRightarrow{\text{Hyp.~(IX.b)}}
\text{Bridge.}
\]
Thread~A uses Cauchy nets;
Thread~B uses compactness;
both originate from
Lemma~\ref{lem:ftb}.
The Bridge theorem is conditional
on Hyp.~(IX.b) (identification problem);
all other results are unconditional.
No eigenfunction analysis;
$\zeta$-connection speculative only.

\textit{Modules:}
(I)~Thread~A
(Secs.~\ref{sec:setup},\ref{sec:strong-limit},\ref{sec:friedrichs});
(II)~Thread~B (Sec.~\ref{sec:mosco});
(III)~Rate (Sec.~\ref{sec:rate},
imports Paper~VIII only);
(IV)~Bridge (Sec.~\ref{sec:bridge},
imports Paper~IX only,
\emph{conditional on Hyp.~(IX.b)});
(V)~Consequences (Sec.~\ref{sec:consequences}).

\textit{Changes v44\,$\to$\,v45:}
Fix~1: lem:monotone uses
$E_{M_\xi}([-c,c])\le E_{M_\xi}([-c',c'])$
via RS~Thm.~VIII.3.
Fix~3: Bridge proof restated as
``same net, unique limit in
Hausdorff space''.
lem:wc unchanged.

\textit{Change v45\,$\to$\,v46:}
Hyp.~(IX.b): identification
$\HSOT=\overline{\Hspectral}$ demoted
from Paper~IX import to explicit
structural hypothesis
(Paper~IX, Open Problem~7;
not proved in Paper~IX).
Bridge theorem now conditional;
all unconditional results unchanged.

\tableofcontents

%%═══════════════════════════════════════════════════════════════════
\section{Setup and Foundational Lemmas}
\label{sec:setup}
%%═══════════════════════════════════════════════════════════════════

The Fourier transform is
$(\Fc f)(\xi):=
\int_{\mathbb{R}}f(x)e^{-i\xi x}dx$
(Plancherel:
$\|\Fc f\|=\sqrt{2\pi}\|f\|$;
\cite{ReedSimon1975}, Ch.~IX).
Set
\begin{equation}\label{eq:ops}
  \Bc:=\Fc^{-1}\chi_{[-c,c]}\Fc,
  \quad
  \Hc:=P^*\mathcal{H}_c^{[0,1]}P,
  \quad
  K_c(x,y)=\tfrac{\sin c(x-y)}{\pi(x-y)},
\end{equation}
where $P:L^2(\mathbb{R})\to L^2([0,1])$
is restriction and $P^*$ extension
by zero.
Then $\Hc$ is self-adjoint,
$0\le\Hc\le I$, $\|\Hc\|\le 1$
(\cite{Slepian1961};
\cite{ReedSimon1975}, Thm.~VI.1).
\hfill\textbf{(F.1)}\label{F1}

\begin{lemma}[Coordinatisation lemma]
\label{lem:ftb}
For all $c>0$, $f\in L^2(\mathbb{R})$:
\begin{equation}\label{eq:ftb}
  \ip{\Hc f}{f}=\|\Bc(P^*Pf)\|^2,
  \qquad
  \mathcal{H}_c^{[0,1]}=P\Bc P^*.
\end{equation}
All subsequent properties of $\Hc$
follow from \eqref{eq:ftb} and
$[-c,c]\subset[-c',c']$.
\end{lemma}

\begin{proof}
Since $(\Fc K_c)(\xi)=\chi_{[-c,c]}(\xi)$
(\cite{Slepian1961}, Eq.~(3)),
the convolution theorem gives
$K_c*h=\Bc h$
(\cite{ReedSimon1975}, Thm.~IX.2).
For $g\in L^2([0,1])$, extending by
zero and restricting:
$(K_c*P^*g)(x)
=\int_0^1 K_c(x-y)g(y)\,dy$,
so $\mathcal{H}_c^{[0,1]}=P\Bc P^*$.
Then $\ip{\Hc f}{f}
=\ip{\Bc P^*Pf}{P^*Pf}
=\|\Bc(P^*Pf)\|^2$.
\end{proof}

\begin{lemma}[Monotonicity ---
  Thread~B]
\label{lem:monotone}
$c\le c'\Rightarrow\ip{\Hc f}{f}
\le\ip{\mathcal{H}_{c'}f}{f}$.
\end{lemma}
\begin{proof}
$\Bc = E_{M_\xi}([-c,c])$
is the spectral projection of the
multiplication operator $M_\xi$
associated to $[-c,c]$.
Since $[-c,c]\subset[-c',c']$,
the spectral measures satisfy
$E_{M_\xi}([-c,c])\le
E_{M_\xi}([-c',c'])$
as orthogonal projections
(\cite{ReedSimon1975}, Thm.~VIII.3),
i.e.\ $\Bc\le B_{c'}$;
apply \eqref{eq:ftb}.
\end{proof}

\begin{lemma}[HS-compactness ---
  Thread~B]
\label{lem:hs-compact}
$\widetilde{\mathcal{H}}_{c_0}$
is compact for each fixed $c_0>0$.
\end{lemma}
\begin{proof}
$\|\mathcal{H}_{c_0}^{[0,1]}\|_{\mathrm{HS}}^2
\le c_0/\pi$
(Fubini $+$ Parseval);
\cite{ReedSimon1975}, Thm.~VI.22.
\end{proof}

\begin{lemma}[Dense domain ---
  Thread~A]
\label{lem:dense}
$\mathcal{D}:=\{f:\Hc f\text{ converges}
\text{ in }L^2\}\supset C_c^\infty$.
\end{lemma}
\begin{proof}
For $f\in C_c^\infty$ and
$h:=P^*Pf\in L^2(\mathbb{R})$:
\[
  \|\Hc f - P^*Pf\|
  \le \|\Bc h - h\|
  = \|(\Fc^{-1}\chi_{[-c,c]}\Fc - I)h\|
  \to 0
\]
as $c\to\infty$, by strong convergence
of the spectral projections
$\chi_{[-c,c]}(M_\xi)$ to the identity
as $[-c,c]\nearrow\mathbb{R}$
(\cite{ReedSimon1975}, Thm.~VIII.5
and the spectral theorem;
$\|P^*P\|\le 1$).
\end{proof}

\begin{remark}
\label{rem:dense-trivial}
Once $\Hc\xrightarrow{\mathrm{SOT}}\Hstrong$,
one has $\mathcal{D}=L^2$;
Lemma~\ref{lem:dense} is
Thread~A pre-input only.
\end{remark}

%%═══════════════════════════════════════════════════════════════════
\section{Thread~A: Strong Operator Limit}
\label{sec:strong-limit}
%%═══════════════════════════════════════════════════════════════════

\begin{lemma}[SOT extension principle]
\label{lem:sc}
Let $(A_c)$ be a directed net in
$\mathcal{B}(\mathcal{H})$ with
$\sup_c\|A_c\|\le M<\infty$.
If $\lim_c A_cg$ exists for all $g$
in a dense $\mathcal{D}$,
then $\exists!\,A\in\mathcal{B}$
with $\|A\|\le M$ and
$A_c\xrightarrow{\mathrm{SOT}}A$.
\end{lemma}
\begin{proof}
For $f$, $g\in\mathcal{D}$ with
$\|f-g\|<\varepsilon/(3M)$:
$\|A_cf-A_{c'}f\|
\le 2M\|f-g\|+\|A_cg-A_{c'}g\|
<\varepsilon$
for $c,c'>c_0(g,\varepsilon)$.
Hence $(A_cf)$ Cauchy;
$Af:=\lim_cA_cf$ is bounded,
linear, unique
(\cite{ReedSimon1975}, Thm.~I.26).
\end{proof}

\begin{lemma}
\label{lem:hlim-bounded}
$\exists!\,\Hstrong\in\mathcal{B}(L^2)$,
$\|\Hstrong\|\le 1$,
with $\Hc\xrightarrow{\mathrm{SOT}}\Hstrong$.
\end{lemma}
\begin{proof}
Lem.~\ref{lem:sc}; $M=1$;
$\mathcal{D}$ from Lem.~\ref{lem:dense}.
\end{proof}

\begin{corollary}\label{cor:qlim-op}
$\qlim(f,f)=\ip{\Hstrong f}{f}$;
$\Hstrong=\Hstrong^*\ge 0$;
$\sigma(\Hstrong)\subset[0,1]$.
\end{corollary}

%%═══════════════════════════════════════════════════════════════════
\section{Thread~B: Mosco Convergence}
\label{sec:mosco}
%%═══════════════════════════════════════════════════════════════════

\begin{definition}[Limit form]
\label{def:qlim}
The function $c\mapsto\ip{\Hc f}{f}$
is bounded above (by $\|f\|^2$) and
monotonically increasing
on $(0,\infty)$
(Lem.~\ref{lem:monotone}),
hence converges in $\mathbb{R}$;
set $\qlim(f,f):=\lim_{c\to\infty}
\ip{\Hc f}{f}=\sup_{c>0}\ip{\Hc f}{f}$
and extend by polarisation.
\end{definition}

\begin{definition}[Mosco convergence
  \cite{Mosco1969}]
\label{def:mosco}
$(\qc)\xrightarrow{\mathrm{M}}\qlim$
iff (M1) $f_n\rightharpoonup f$
$\Rightarrow$
$\qlim(f,f)\le\liminf_n q_{c_n}(f_n,f_n)$;
and (M2) $\forall f$,
$\exists f_n\to f$:
$q_{c_n}(f_n,f_n)\to\qlim(f,f)$.
\end{definition}

\begin{lemma}[Form continuity]
\label{lem:wc}
Compact $T$, $f_n\rightharpoonup f$
$\Rightarrow$
$\ip{Tf_n}{f_n}\to\ip{Tf}{f}$.
\end{lemma}
\begin{proof}
$Tf_n\to Tf$ strongly
(\cite{ReedSimon1975}, Thm.~VI.1);
weakly convergent sequences are
bounded (Thm.~II.12);
hence strong--weak continuity of
$\ip{\cdot}{\cdot}$.
\end{proof}

\begin{theorem}[Mosco convergence]
\label{thm:mosco}
$\qc\xrightarrow{\mathrm{M}}\qlim$.
\end{theorem}
\begin{proof}
(M2) is immediate from
Def.~\ref{def:qlim}.
For (M1): given $f_n\rightharpoonup f$
and $c_0>0$, monotonicity gives
$q_{c_n}(f_n,f_n)\ge q_{c_0}(f_n,f_n)$
for $c_n\ge c_0$;
Lems.~\ref{lem:hs-compact}
and~\ref{lem:wc} give
$q_{c_0}(f_n,f_n)\to q_{c_0}(f,f)$;
take $\sup_{c_0}$.
\end{proof}

%%═══════════════════════════════════════════════════════════════════
\section{Identification and Closedness}
\label{sec:friedrichs}
%%═══════════════════════════════════════════════════════════════════

\begin{theorem}[Closedness and
  Riesz--Kato identification]
\label{thm:friedrichs}
$\qlim$ is closed and
$T_{\qlim}=\Hstrong$.
\end{theorem}
\begin{proof}
Closedness: bounded symmetric
forms are closed
(\cite{ReedSimon1975}, Thm.~VIII.15).
Identification:
$\ip{\Hstrong f}{g}=\qlim(f,g)$
(Cor.~\ref{cor:qlim-op});
uniqueness of Riesz representative.
\end{proof}

\begin{remark}
Thread~A produced $\Hstrong$;
Thread~B produced $\qlim$;
Thm.~\ref{thm:friedrichs}
identifies them.
\end{remark}

%%═══════════════════════════════════════════════════════════════════
\section{Quantitative Rate
  (External, Paper~VIII)}
\label{sec:rate}
%%═══════════════════════════════════════════════════════════════════

\begin{theorem}[Quantitative rate]
\label{thm:rate-ext}
$0\le\qlim(f,f)-\qc(f,f)
\le C_\kappa c^{-1/4}\|f\|^2$,
$c\ge c_0(\kappa)$.
(External: Paper~VIII, Cor.~4.3;
no other section depends on this
theorem.)
\end{theorem}

\begin{remark}[Obstruction]
\label{rem:no-indep-proof}
An internal proof requires uniform
eigenvector asymptotics;
mixing fixed-$c$ eigenbasis with
$c\to\infty$ weights is invalid
(Problem~\ref{prob:rate}).
\end{remark}

%%═══════════════════════════════════════════════════════════════════
\section{The Bridge Theorem}
\label{sec:bridge}
%%═══════════════════════════════════════════════════════════════════

Sole import point for Paper~IX.
Assume (Paper~IX, Thms.~2.5, 3.2,
Rem.~2.6):
\begin{enumerate}[(IX.a)]
\item\label{IXa}
  $\Hspectral:=
  \sum_n\lambda_n^{(\infty)}
  \ip{\cdot}{\Phi_n^{(\infty)}}
  \Phi_n^{(\infty)}$
  is densely defined, symmetric,
  closable, with
  $\|\Hspectral f\|\le\|f\|$;
  hence
  $\overline{\Hspectral}\in
  \mathcal{B}(L^2)$.
  \emph{(Unconditional in Paper~IX,
  Thm.~3.3 and Rem.~1.3.)}
\item\label{IXb}
  \emph{(Hypothesis; identification
  problem from Paper~IX,
  Open Problem~7;
  not proved in Paper~IX.)}
  The net $(\Hc)_{c\in(0,\infty)}$
  converges in the strong operator
  topology to a limit $\HSOT$ which
  coincides with the closure of the
  spectral operator:
  \[
    \Hc\xrightarrow{\mathrm{SOT}}\HSOT
    \quad\text{and}\quad
    \HSOT=\overline{\Hspectral}.
  \]
  The existence of the SOT-limit is
  unconditional
  (Lem.~\ref{lem:hlim-bounded});
  the identification
  $\HSOT=\overline{\Hspectral}$
  is a structural hypothesis
  not yet established in the series.
\end{enumerate}

\begin{theorem}[Bridge theorem
  (conditional on Hyp.~\ref{IXb})]
\label{thm:bridge}
\emph{Conditional on
Hyp.~\hyperref[IXb]{(IX.b)}.}
$\Hstrong=\overline{\Hspectral}$.
\end{theorem}
\begin{proof}
The net $(\Hc)_{c\in(0,\infty)}$
ordered by $c\le c'$ converges
in the SOT topology on
$\mathcal{B}(L^2)$ to both
$\Hstrong$ (Lem.~\ref{lem:hlim-bounded})
and $\HSOT=\overline{\Hspectral}$
(Hyp.~\hyperref[IXb]{(IX.b)}).
Limits of a convergent net in a
Hausdorff topological space are
unique; hence
$\Hstrong=\overline{\Hspectral}$.
\end{proof}

\begin{remark}[Separation of existence
  and identification]
\label{rem:bridge-sep}
The SOT-limit $\Hstrong$ exists
unconditionally
(Lem.~\ref{lem:hlim-bounded})
and satisfies $T_{\qlim}=\Hstrong$
(Thm.~\ref{thm:friedrichs}).
The identification
$\Hstrong=\overline{\Hspectral}$
requires Hyp.~\hyperref[IXb]{(IX.b)}
and is the content of
Thm.~\ref{thm:bridge}.
The two questions ---
\emph{existence} of the SOT-limit vs.
\emph{identification} with the
spectral closure --- are logically
independent; conflating them was
the source of the earlier circularity.
\end{remark}

\begin{remark}
No eigenfunction analysis:
under Hyp.~\hyperref[IXb]{(IX.b)},
equality follows from uniqueness of
the SOT-limit in the Hausdorff
topology on $\mathcal{B}(L^2)$.
$\Hlim:=\Hstrong=\overline{\Hspectral}$
(conditional).
\end{remark}

%%═══════════════════════════════════════════════════════════════════
\section{Spectral Consequences}
\label{sec:consequences}
%%═══════════════════════════════════════════════════════════════════

\begin{corollary}[Resolvent convergence]
\label{cor:resolvent}
For $z\notin[0,1]$:
$(\Hc-z)^{-1}
\xrightarrow{\mathrm{SOT}}
(\Hlim-z)^{-1}$.
\end{corollary}
\begin{proof}
Thm.~\ref{thm:mosco};
$\sigma(\Hlim)\subset[0,1]$;
\cite{KuwaeShioya2003}, Thm.~2.4.
\end{proof}

\begin{remark}[Unconditional summary]
$\Hlim^*=\Hlim\ge 0$;
$T_{\qlim}=\Hlim$;
$\qc\xrightarrow{\mathrm{M}}\qlim$;
$(\Hc-z)^{-1}\xrightarrow{\mathrm{SOT}}
(\Hlim-z)^{-1}$;
rate $O(c^{-1/4})$ (Paper~VIII).
The identification
$\Hstrong=\overline{\Hspectral}$
is conditional on
Hyp.~\hyperref[IXb]{(IX.b)}.
\end{remark}

\begin{remark}[Speculative:
  Riemann zeros]
\label{rem:rh}
\emph{No unconditional claim.}
Under (H1) $\{\Phi_n^{(\infty)}\}$
complete, $\Hlim$ ess.\ self-adjoint
(Paper~IX, Hyp.~4.1);
and (H2) zeros $\rho$ in
$\sigma_{\mathrm{app}}(\Hlim)$
(\cite{CCM2025}, Thm.~1.3):
$\rho\in\sigma(\Hlim)\subset\mathbb{R}$,
hence $\mathrm{Re}(\rho)=1/2$.
(H2) has no proved operator
identification.
\end{remark}

%%═══════════════════════════════════════════════════════════════════
\section{Open Problems}
\label{sec:problems}
%%═══════════════════════════════════════════════════════════════════

\begin{problem}[Completeness]
\label{prob:completeness}
Prove $\{\Phi_n^{(\infty)}\}$
complete in $L^2(\mathbb{R})$.
\end{problem}

\begin{problem}[Identification]
\label{prob:identification}
Prove Hyp.~\hyperref[IXb]{(IX.b)}:
$\HSOT=\overline{\Hspectral}$.
This is Paper~IX, Open Problem~7,
and the central open problem
for the Bridge theorem.
\end{problem}

\begin{problem}[Independent rate]
\label{prob:rate}
Prove rate $O(c^{-1/4})$ without
Paper~VIII.
\end{problem}

\begin{problem}[Quantitative
  resolvent]
\label{prob:resolvent}
Find $F(c,z)\to 0$ with
$\|(\Hc-z)^{-1}-(\Hlim-z)^{-1}\|
\le F(c,z)$.
\end{problem}

\begin{problem}[Spectrum]
\label{prob:spectrum}
Is $\sigma(\Hlim)=[0,1]$?
\end{problem}

%%%─── Bibliography
\begin{thebibliography}{99}

\bibitem{PaperV}
[Author],
\textit{Effective bandwidth and
$L^2$-concentration of PSWFs},
Paper~V, preprint 2025.

\bibitem{PaperVIII}
[Author],
\textit{Non-cancellation, sharp
$c^{-1/2}$ lower bound,
and operator-norm asymptotics},
Paper~VIII, preprint 2026.

\bibitem{PaperIX}
[Author],
\textit{Domain, deficiency indices,
self-adjointness},
Paper~IX, preprint 2026.

\bibitem{Kato1966}
T.~Kato,
\textit{Perturbation Theory for
Linear Operators},
Springer, 1966.

\bibitem{ReedSimon1975}
M.~Reed and B.~Simon,
\textit{Methods of Modern
Mathematical Physics, Vols.~I--II},
Academic Press, 1972--1975.

\bibitem{Mosco1969}
U.~Mosco,
\textit{Convergence of convex sets
and of solutions of variational
inequalities},
Adv.\ Math.\ \textbf{3} (1969),
510--585.

\bibitem{KuwaeShioya2003}
K.~Kuwae and T.~Shioya,
\textit{Convergence of spectral
structures},
Comm.\ Anal.\ Geom.\ \textbf{11}
(2003), 599--673.

\bibitem{Slepian1961}
D.~Slepian and H.~O.~Pollak,
\textit{Prolate spheroidal wave
functions, Fourier analysis
and uncertainty---I},
Bell Syst.\ Tech.\ J.\
\textbf{40} (1961), 43--63.

\bibitem{CCM2025}
A.~Connes, C.~Consani,
and H.~Moscovici,
\textit{Zeta Spectral Triples},
arXiv:2511.22755, 2025.

\bibitem{Rudin1991}
W.~Rudin,
\textit{Functional Analysis},
2nd ed., McGraw-Hill, 1991.

\bibitem{Conway1990}
J.~B.~Conway,
\textit{A Course in Functional
Analysis},
2nd ed., Springer, 1990.

\end{thebibliography}

\end{document}
